% Section file: Twisted Holography and Quantum Gravity -- Non-perturbative completion
%
% Develops the D3-instanton sector of the twisted-holographic
% partition function, the Borcherds product packaging the perturbative
% Heegner tower, and the conditional resurgent assembly with instanton
% and Schwinger tails.

\section{Non-perturbative completion of the twisted-holographic partition function}
\label{sec:thqg-nonperturbative}
\label{ch:thqg-nonperturbative}
\index{non-perturbative completion!twisted holography|textbf}
\index{D3-instanton!twisted-holographic|textbf}
\index{Borcherds product!non-perturbative resummation}

% =====================================
% Local macros
% =====================================
\providecommand{\HS}{\operatorname{HS}}
\providecommand{\Tr}{\operatorname{Tr}}
\providecommand{\Det}{\operatorname{Det}}
\providecommand{\vol}{\operatorname{vol}}
\providecommand{\Res}{\operatorname{Res}}
\providecommand{\rig}{\mathrm{rig}}

Perturbative finiteness of the twisted-holographic partition function
(Theorem~\ref{thm:thqg-perturbative-finiteness}) secures convergence
of the genus expansion on the Costello HT locus. Away from that locus,
the Heegner-patterned BV obstruction tower
$\{c_n^{\mathrm{triangle}}\}_{n\geq 1}$ of
Theorem~\ref{thm:htbb-heegner-all-order} grows as
$|c_n^{\mathrm{triangle}}|\sim\exp(\pi\sqrt n)$, faster than any
polynomial in~$n$ but subfactorial. This coefficient estimate alone
does not imply a Gevrey-$1$ asymptotic series, Borel summability, or
Borel-plane singularities. The present section records the
Borcherds infinite product for~$\Phi_{10}^{\mathrm{un}}$ as the automorphic package
of the Heegner tower and isolates the additional resurgent hypotheses
under which Euclidean D$3$-branes and Schwinger sectors assemble a
non-perturbative bulk--boundary--line coupling.

The convention throughout this section is the non-OP Igusa square
\[
 \Phi_{10}^{\mathrm{un}}=\Delta_5^2,
\]
where \(\Delta_5\) is the primitive weight-\(5\) BKM denominator.  The
monic product is \(D_5=64^{-1}\Delta_5\), and the
Oberdieck--Pixton reduced-DT scalar convention would be
\[
 \Phi_{10}^{\mathrm{OP}}=D_5^2=4096^{-1}\Delta_5^2,\qquad
 Z_{\mathrm{OP/DT}}=-(\Phi_{10}^{\mathrm{OP}})^{-1}
 =-D_5^{-2}=-4096\,\Delta_5^{-2}.
\]
No reduced-DT normalisation is used below unless explicitly named.

The section develops this completion in five subsections:
(\S\ref{subsec:thqg-np-d3-instanton}) the D$3$-instanton amplitude and
its conditional Borel-plane interpretation;
(\S\ref{subsec:thqg-np-borcherds-resummation}) the Borcherds product
as automorphic packaging of the Heegner tower;
(\S\ref{subsec:thqg-np-schwinger-boundary}) Schwinger pair creation at
the Humbert boundary;
(\S\ref{subsec:thqg-np-climax-theorem}) the climax non-perturbative
completion theorem;
(\S\ref{subsec:thqg-np-multipath}) multi-path verification of the
non-perturbative completion.

% ======================================================================
\subsection{The D\texorpdfstring{$3$}{3}-instanton amplitude on rigid K3 curves}
\label{subsec:thqg-np-d3-instanton}

\begin{definition}[Rigid K3 curve class]
\label{def:thqg-np-rigid-curve}
\index{rigid curve!K3}
A holomorphic curve $C\subset K3$ is \emph{rigid} if
$\mathcal{N}_{C/K3} = \mathcal{O}_{\bP^1}(-2)$, equivalently
$[C]^2 = -2$ in the Mukai lattice $\Lambda_{K3} = 2U\oplus 3(-E_8)$
of signature $(3, 19)$. Denote by
$H_2^+(K3)_{\rig} := \{[C]\in H_2(K3,\bZ)\mid [C]^2 = -2,\,C\in|[C]|\text{ rigid}\}$
the set of rigid curve classes modulo $\mathrm{Aut}(K3)$-orbits;
$H_2^+(K3)_{\rig}$ is finite at each complex-moduli point
$\tau\in\overline{\mathcal{A}_2}$ by Nikulin's classification of
$(-2)$-root systems
in~$\Lambda_{K3}$~\cite[Thm.~1.12.2]{Nikulin1980}.
\end{definition}

\begin{proposition}[Euclidean D$3$-instanton amplitude;
\ClaimStatusProvedHere]
\label{prop:thqg-np-d3-amplitude}
\index{D3-instanton!amplitude formula|textbf}
\index{Vafa F-theory!D3 on K3 curve}
In the Costello--Li holomorphic-topological twist of type~IIB on
$\bR^3\times K3\times\bC^2$, dual by $F$-theory on
$K3\times K3$~\cite[\S 2]{VafaFTheory1996} to M-theory on
$\bR^3\times K3\times\bC^2$, a Euclidean D$3$-brane wrapping a rigid
curve $C$ with $[C]\in H_2^+(K3)_{\rig}$ contributes a
non-perturbative amplitude
\begin{equation}\label{eq:thqg-np-d3-amp}
 \mathcal{A}^{\mathrm{D3}}_{[C]}(\hbar)
 \;=\; \mathcal{Z}^{(1)}_{\det}(C)\cdot e^{-2\pi\vol(C)/\hbar},
 \qquad \vol(C) \;=\; \int_C\omega_{K3},
\end{equation}
where $\omega_{K3}$ is the K\"ahler form of the complex-moduli point
$\tau\in\overline{\mathcal{A}_2}$ and
$\mathcal{Z}^{(1)}_{\det}(C) = \Det^{-1}(-\bar\partial_{C,\mathcal{N}})$
is the Quillen--Bismut determinant~\cite[\S 7]{BismutGilletSoule1988}
of the $\bar\partial$-operator on the normal bundle
$\mathcal{N}_{C/K3} = \mathcal{O}_{\bP^1}(-2)$. The amplitude is
identically zero order by order in the formal $\hbar$-expansion.
\end{proposition}

\begin{proof}
Rigidity $[C]^2 = -2$ implies the deformation-moduli space of $C$ is
a point: $H^0(C,\mathcal{N}_{C/K3}) =
H^0(\bP^1,\mathcal{O}_{\bP^1}(-2)) = 0$ by Serre vanishing, and
$H^1(C,\mathcal{N}_{C/K3}) =
H^1(\bP^1,\mathcal{O}_{\bP^1}(-2)) = \bC$ (one obstruction mode,
absorbed by the BV antifield pairing). The classical action is
$S_{\mathrm{cl}}(C) = 2\pi\int_C\omega_{K3} = 2\pi\vol(C)$ by the
K\"ahler calibration for holomorphic cycles. The one-loop fluctuation
determinant is the Quillen--Bismut determinant
$\Det^{-1}(-\bar\partial_{C,\mathcal{N}})$ of the normal-bundle
$\bar\partial$-operator, finite by~\cite[Thm.~7.1]{BismutGilletSoule1988}
for~$\mathcal{N} = \mathcal{O}(-2)$ on $\bP^1$. The amplitude
\eqref{eq:thqg-np-d3-amp} follows from the semiclassical BV
evaluation of the D$3$ path integral at the rigid classical solution.

The factor $e^{-2\pi\vol(C)/\hbar}$ vanishes identically in every
formal power of $\hbar$: for all $n\geq 0$, the Taylor coefficient of
$e^{-2\pi\vol(C)/\hbar}$ at $\hbar = 0$ in the formal variable
$\hbar$ is zero. D$3$-instantons are therefore invisible in the
perturbative BV obstruction tower but contribute beyond all orders
through the non-perturbative Borel-plane structure.
\end{proof}

\begin{remark}[Borel-plane location of the D$3$-instanton]
\label{rem:thqg-np-d3-borel-plane}
\index{Borel plane!D3-instanton singularity}
The Borel transform of a Gevrey-$1$ asymptotic series
$f(\hbar) = \sum_{n\geq 0}f_n\hbar^n$ with $|f_n|\leq C^n n!$ is the
analytic function $\mathcal{B}[f](\zeta) := \sum_{n\geq 0}f_n\zeta^n/n!$,
a germ of an analytic function near $\zeta = 0$. Its singularities
in the Borel plane $\zeta\in\bC$ occur at instanton actions only
after a resurgent transseries structure has been supplied
(\`Ecalle--Voros~\cite{EcalleVoros1981}). Thus the D$3$-brane on a
rigid curve $C$ gives the candidate action
$\zeta = 2\pi\vol(C)$ and candidate one-loop coefficient
$\mathcal{Z}^{(1)}_{\det}(C)$. To turn this candidate into a theorem
one must prove that the twisted-holographic $\hbar$-series is
Gevrey-$1$, admits sectorial analytic continuation of its Borel
transform, and has the D$3$ saddle as an actual singularity with that
residue. The Heegner coefficient growth used above does not supply
these analytic hypotheses.
\end{remark}

% ======================================================================
\subsection{The Borcherds product as non-perturbative resummation}
\label{subsec:thqg-np-borcherds-resummation}

\begin{theorem}[Borcherds product of $\Delta_5$ packages the
Heegner tower; \ClaimStatusProvedHere]
\label{thm:thqg-np-borcherds-resummation}
\index{Borcherds product!resummation of Heegner tower|textbf}
\index{Igusa cusp form!infinite product}
The perturbative Heegner tower $\sum_{n\geq 1}c_n^{\mathrm{triangle}}\hbar^n$
of Theorem~\ref{thm:htbb-heegner-all-order}, interpreted as the
$\hbar$-coupling tower of the twisted-holographic bulk--boundary--line
triangle, has scalar automorphic package given by the
Gritsenko--Nikulin product for the primitive denominator $\Delta_5$.
Write
\[
 \phi_{0,1}(\tau,z)=\sum_{n\geq 0,\ r\in\mathbb Z} f(n,r)q^n y^r.
\]
The monic Borcherds product and the unnormalised Igusa square are
\begin{equation}\label{eq:thqg-np-borcherds-product}
 D_5(Z)
 \;=\;
 p^{1/2}q^{1/2}y^{1/2}
 \prod_{(m, n, r)>0}\!(1 - p^m q^n y^r)^{f(mn,r)}
 \;=\;64^{-1}\Delta_5(Z),
\end{equation}
and hence
\[
 \Phi_{10}^{\mathrm{un}}(Z)
 =\Delta_5(Z)^2
 =4096\,p q y
 \prod_{(m,n,r)>0}(1-p^m q^n y^r)^{2f(mn,r)}.
\]
The integers $f(mn,r)$ are signed denominator exponents.  Ordinary
D$3$-sector multiplicities, BPS Hilbert-space dimensions, or Hall
source dimensions are not consequences of this scalar identity alone.
\end{theorem}

\begin{proof}
By Borcherds~$1998$ Theorem~$10.1$~\cite[Thm.~10.1]{Borcherds1998},
the primitive denominator $\Delta_5$ is the singular-theta lift of the
weight-zero index-one weak Jacobi form $\phi_{0,1}$; in the
Gritsenko--Nikulin normalisation the monic product is
\eqref{eq:thqg-np-borcherds-product}. Taking $\log$,
\[
 \log D_5(Z) \;=\; \tfrac12\log(pqy)
 \;+\; \sum_{(m, n, r)>0}f(mn,r)\log(1 - p^m q^n y^r).
\]
Expanding each $\log(1 - p^m q^n y^r) = -\sum_{k\geq 1}(p^m q^n y^r)^k/k$
and collecting the $p^n$-Fourier coefficient yields (after the
$\eta^{24}$ normalisation absorbs the one-loop paramodular anomaly,
as in Theorem~\ref{thm:htbb-heegner-all-order})
$f(n,\ell) + (\text{denominator-identity lower orders})$, the
perturbative $\hbar^n$-order
of~Theorem~\ref{thm:htbb-heegner-all-order}.

Squaring gives $\Phi_{10}^{\mathrm{un}}=\Delta_5^2$ with exponents
$2f(mn,r)$, while the OP branch uses
$\Phi_{10}^{\mathrm{OP}}=D_5^2=4096^{-1}\Delta_5^2$.  This proves the
automorphic product identity and the Fourier--Jacobi matching with the
Heegner tower. It does not prove that the same product is the Borel sum
of a Gevrey-$1$ $\hbar$-series; that is the conditional resurgent
interpretation isolated below.
\end{proof}

\begin{remark}[Signed scalar exponent and Hall recognition]
\label{rem:thqg-np-multiplicity}
\index{D3-instanton!multiplicity}
The coefficient $f(mn,r)$ in \eqref{eq:thqg-np-borcherds-product} is
the signed BKM root supercharacter of the $\Delta_5$ denominator.  A
reading as an ordinary degeneracy of
$\mathbf{H}_{\Delta_5}$-module BPS states with K\"ahler-modulus
decomposition $(m,n,r)$ is available only after the Vol~III finite
Hall--Borcherds recognition datum supplies compact source primitives,
parity blocks, product and coproduct matrices, super-brackets, Hopf
pairing, radical quotient, no-extra-relations equality, PBW comparison,
centre and associator control, and compatible finite-window transitions.
Three scalar verification paths agree before that recognition step:
(i)~the Eichler--Zagier tabulation of $\phi_{0,1}$-coefficients
\cite[\S 9]{EichlerZagier1985}; (ii)~the Harvey--Moore heterotic threshold
integral~\cite[\S 4]{HarveyMoore1996}; (iii)~the
Gritsenko--Nikulin denominator-identity decomposition of
$\Phi_{10}^{\mathrm{un}}$~\cite[Thm.~2.1]{GritsenkoNikulin1998} into BKM-root
contributions.
\end{remark}

% ======================================================================
\subsection{Schwinger pair creation at the Humbert boundary}
\label{subsec:thqg-np-schwinger-boundary}

The Schwinger effect at the Humbert boundary, derived in detail as
Theorem~\ref{thm:htbb-schwinger-humbert} of
\S\ref{sec:htbb-schwinger-humbert}, is reproduced here for the
climax-theorem assembly. The pair-creation rate is
\begin{equation}\label{eq:thqg-np-schwinger-rate}
 \Gamma^{\mathrm{Schwinger}}_{H_1}(\hbar)
 \;=\; \frac{(\vol H_1)^2}{(2\pi\hbar)^2}\,e^{-\pi m^2/|q|\hbar},
\end{equation}
with $m^2$ the $\mathbf{H}_{\Delta_5}$-module mass squared at the
Humbert locus, defined only after the finite Hall--Borcherds
recognition datum of Remark~\ref{rem:thqg-np-multiplicity}, and
$q\in\{1, \ldots, 7\}$ the $\mu_8$-holonomy charge sector. At Lusztig
specialisation $\hbar_{\mathrm{qg}}^2 = -1/8$,
\[
 \Gamma^{\mathrm{Schwinger}}_{H_1}(\hbar_{\mathrm{qg}})
 \;=\; \frac{2(\vol H_1)^2}{\pi^2}\,e^{-2\pi\sqrt 2\,m^2/|q|}.
\]

\begin{remark}[Schwinger corner in the bulk--boundary--line triangle]
\label{rem:thqg-np-schwinger-corner}
\index{HT triangle!Schwinger corner}
The Schwinger pair-creation rate
\eqref{eq:thqg-np-schwinger-rate} contributes to the
twisted-holographic bulk--boundary--line triangle through the
boundary/line corner only after the finite Hall--Borcherds recognition
datum has supplied the $\mathbf{H}_{\Delta_5}$ source category and its
transport functor: the Wilson-line transport
$\mathcal{T}_\gamma\colon \mathbf{H}_{\Delta_5}|_{\gamma(0)}\to
\mathbf{H}_{\Delta_5}|_{\gamma(1)}$ along a boundary-crossing path
$\gamma$ develops a non-unitary defect
$\Gamma^{\mathrm{Schwinger}}_{H_1}(\hbar)\cdot\mathcal{P}_{\mathrm{pair}}$
per Theorem~\ref{thm:htbb-schwinger-humbert}. The defect is invisible
in the perturbative $\hbar$-expansion of the triangle
(Theorem~\ref{thm:htbb-heegner-all-order}) and appears only through
the non-perturbative $e^{-\pi m^2/|q|\hbar}$ tail.
\end{remark}

% ======================================================================
\subsection{The climax non-perturbative completion theorem}
\label{subsec:thqg-np-climax-theorem}

\begin{theorem}[Conditional non-perturbative presentation of the
twisted-holographic partition function; \ClaimStatusConditional]
\label{thm:thqg-np-d3-completion}
\index{THQG partition function!non-perturbative completion|textbf}
\index{twisted holography!Borel resummation}
Let
$\mathcal{Z}^{\mathrm{pert}}_{\mathrm{HT}}(\hbar, \tau)
:= \sum_{n\geq 0}\hbar^n c_n^{\mathrm{triangle}}(\tau)$
denote the perturbative twisted-holographic bulk--boundary--line
coupling of Theorem~\ref{thm:htbb-heegner-all-order}, viewed as a
formal $\hbar$-series with Heegner-Humbert-class coefficients on the
complex-moduli base $\tau\in\overline{\mathcal{A}_2}$. Assume:
\begin{enumerate}[label=\textup{(\alph*)}]
\item the formal $\hbar$-series is Gevrey-$1$ resurgent and its
Borel singularities at D$3$ actions have the residues proposed in
Remark~\ref{rem:thqg-np-d3-borel-plane};
\item the Schwinger tail of Theorem~\ref{thm:htbb-schwinger-humbert}
is the boundary-sector non-perturbative correction in the same
transseries;
\item the formal substitution $p^\hbar$ below identifies the
Fourier-Jacobi expansion with the chosen $\hbar$-expansion.
\item any $\mathbf H_{\Delta_5}$ module, transport, or BPS
state-space interpretation is made only after the finite
Hall--Borcherds recognition datum of
Remark~\ref{rem:thqg-np-multiplicity}.
\end{enumerate}
Then the non-perturbative presentation assembling the D$3$-instanton tail
(Proposition~\ref{prop:thqg-np-d3-amplitude}) and the Schwinger
pair-creation tail
(Theorem~\ref{thm:htbb-schwinger-humbert}) is
\begin{equation}\label{eq:thqg-np-completion}
 \mathcal{Z}^{\mathrm{full}}_{\mathrm{HT}}(\hbar, \tau)
 \;=\; \mathcal{Z}^{\mathrm{pert}}_{\mathrm{HT}}(\hbar, \tau)
 \;+\;\!\!\sum_{[C]\in H_2^+(K3)_{\rig}}\!\!
 e^{-2\pi\vol(C)/|\hbar|}\cdot c^{\mathrm{D3}}_{[C]}
 \;+\; \sum_{q = 1}^{7}\Gamma^{\mathrm{Schwinger}}_{H_1, q}(\hbar),
\end{equation}
and~\eqref{eq:thqg-np-completion} equals the Borcherds-product form
\eqref{eq:thqg-np-borcherds-product}, squared and multiplied by
$4096$, for $\Phi_{10}^{\mathrm{un}}(Z)=\Delta_5(Z)^2$ up to the
factor~$\eta^{24}$ absorbing the one-loop paramodular anomaly:
\begin{equation}\label{eq:thqg-np-phi10-identity}
 \mathcal{Z}^{\mathrm{full}}_{\mathrm{HT}}(\hbar, \tau)
 \;=\; \log\!\bigl(\Phi_{10}^{\mathrm{un}}(Z)/\eta(\tau)^{24}\bigr)\bigr|_{p^\hbar},
\end{equation}
with the formal substitution $p^\hbar := \sum_{n\geq 0}p^n\hbar^n/n!$
extending the formal $p$-expansion to a formal $\hbar$-expansion by
Mellin transform. At Lusztig specialisation $\hbar_{\mathrm{qg}}^2 =
-1/8$, all three tails are simultaneously convergent: the
perturbative tail converges by $\mu_8$-banding truncation
(Theorem~\ref{thm:htbb-W18-banding-wilson-line}); the D$3$-instanton
tail converges by
$e^{-4\pi\sqrt 2\vol(C)}\leq e^{-4\pi\sqrt 2}\approx 1.72\times 10^{-8}$
(minimal rigid curve $\vol(C) = 1$); the Schwinger tail converges by
$e^{-2\pi\sqrt 2 m^2/|q|}$ suppression.
\end{theorem}

\begin{proof}
The identity~\eqref{eq:thqg-np-completion} assembles three inputs.

\emph{(i)} The perturbative tail
$\mathcal{Z}^{\mathrm{pert}}_{\mathrm{HT}}$ equals
$\log(\Phi_{10}^{\mathrm{un}}/\eta^{24})$ expanded as a formal power series
in~$p$: this is the content of
Theorem~\ref{thm:htbb-heegner-all-order} combined with Borcherds'
singular-theta lift~\cite[Thm.~10.1]{Borcherds1998}.

\emph{(ii)} The D$3$-instanton tail
$\sum_{[C]}e^{-2\pi\vol(C)/|\hbar|}\cdot c^{\mathrm{D3}}_{[C]}$
is read from the same scalar exponent table as the
$\Delta_5$ product of Theorem~\ref{thm:thqg-np-borcherds-resummation}.
The equality with a physical instanton sum is part of the resurgent and
finite-recognition hypotheses, not a consequence of the scalar product.

\emph{(iii)} The Schwinger pair-creation tail
$\sum_q\Gamma^{\mathrm{Schwinger}}_{H_1, q}(\hbar)$ enters through
the $\mu_8$-gerbe holonomy sectors of the $\Phi_{10}^{\mathrm{un}}$ denominator
identity (Gritsenko--Nikulin~\cite[\S 3]{GritsenkoNikulin1998}): the
seven non-trivial $\mu_8$-holonomy charges index the seven non-trivial
Gritsenko--Nikulin denominator decompositions modulo $\mu_8$-action on
the Humbert boundary.

Combining (i)--(iii) under the four hypotheses in the theorem, the
Borcherds infinite product gives the same transseries presentation as
the perturbative $\log$-expansion plus the D$3$-instanton tail plus
the Schwinger $\mu_8$-tail. This is~\eqref{eq:thqg-np-completion}
with the identification~\eqref{eq:thqg-np-phi10-identity}.

\emph{Lusztig convergence.} At $\hbar^2_{\mathrm{qg}} = -1/8$,
$|\hbar_{\mathrm{qg}}| = 1/(2\sqrt 2)$ and the three suppression
factors are $\mu_8$-banding truncation
(Theorem~\ref{thm:htbb-W18-banding-wilson-line}),
$e^{-4\pi\sqrt 2\vol(C)}$ (D$3$), and $e^{-2\pi\sqrt 2 m^2/|q|}$
(Schwinger), all non-perturbatively suppressed.
\end{proof}

\begin{remark}[Conditional interpretation: Borcherds product as a
non-perturbative moduli space]
\label{rem:thqg-np-moduli-interpretation}
\index{Borcherds product!moduli interpretation}
Under the resurgent hypotheses of
Theorem~\ref{thm:thqg-np-d3-completion} and the finite recognition
hypothesis in Remark~\ref{rem:thqg-np-multiplicity}, the infinite
product \eqref{eq:thqg-np-borcherds-product} has a physical
interpretation as a sum over D$3$-instanton sectors of the moduli space
of twisted $11$D~SUGRA. Each Borcherds factor
$(1 - p^m q^n y^r)^{f(mn,r)}$ carries the signed scalar exponent of
one K\"ahler modulus $(m,n,r)$; an ordinary D$3$-sector multiplicity is
the corresponding parity-refined source dimension only after
finite Hall--Borcherds recognition. The partition function is then read
as the product of per-sector contributions, reflecting the factorisation
of the D$3$-instanton path integral into
independent-modulus sectors. The zero of $\Phi_{10}^{\mathrm{un}}$ at the diagonal
Humbert divisor $\mathrm{diag}\subset\overline{\mathcal{A}_2}$
corresponds to the tachyonic mode of the $(-2)$-root
$[C_{0, 0, -1}]$, whose instanton action vanishes at
$\tau_1 = \tau_2$. This is the gauge-theory mechanism of the
conformal-point singularity on
$\overline{\mathcal{A}_2}$~\cite[\S 4]{GritsenkoNikulin1997}.
\end{remark}

% ======================================================================
\subsection{Multi-path verification of the non-perturbative completion}
\label{subsec:thqg-np-multipath}

\begin{remark}[Five compatibility checks]
\label{rem:thqg-np-multipath}
\index{multi-path verification!non-perturbative completion}
The automorphic product in Theorem~\ref{thm:thqg-np-borcherds-resummation}
and the conditional presentation in
Theorem~\ref{thm:thqg-np-d3-completion} have five compatibility
checks.
\begin{enumerate}[leftmargin=1.8em,label=\textup{(\alph*)}]
\item \emph{Heterotic path.} Harvey--Moore~$1996$
\cite[\S 4]{HarveyMoore1996} evaluate the heterotic one-loop
threshold correction on $K3\times T^2$ as $\log\Phi_{10}^{\mathrm{un}}/\eta^{24}$
directly, with D$3$-instanton tail given by Fourier coefficients of
$\phi_{0,1}$ term by term.
\item \emph{$F$-theory path.} Vafa~$1996$
\cite[\S 2]{VafaFTheory1996} identifies Euclidean D$3$-branes wrapping
K$3$ rigid curves as the $F$-theory instanton tail; the action
$S_{\mathrm{D3}} = 2\pi\vol(C)$ is the K\"ahler calibration.
\item \emph{DMVV path.} The symmetric-product orbifold elliptic genus
of~$K3^{[n]}$ computes the scalar square exponents via the
Dijkgraaf--Moore--Verlinde--Verlinde
formula~\cite{DMVV1997} order by order; agreement term by term with
$2f(mn,r)$ in the $\Phi_{10}^{\mathrm{un}}=\Delta_5^2$ lane.
\item \emph{M$5$-brane instanton path.} Witten's M$5$-instanton
count~\cite[\S 5]{Witten1996M5} on the Kodaira-$I_1$ fibre of the
elliptic K$3$ gives D$3$-instanton contributions at $n$-instanton
level with Kodaira-$I_1$ boundary, matching the Borcherds exponents
directly.
\item \emph{BKM denominator path.} Gritsenko--Nikulin~$1998$
\cite[Thm.~2.1]{GritsenkoNikulin1998} decompose $\Phi_{10}^{\mathrm{un}}$ into BKM
root-multiplicity factors; each factor is one rigid curve class by
Theorem~\ref{thm:thqg-np-borcherds-resummation}, verifying the
signed denominator character of the BKM algebra $\mathfrak{g}_{\Delta_5}$.
\end{enumerate}
All five checks agree on the Fourier coefficients of $\phi_{0,1}$; the
leading values
\[
 f(0,\pm1)=1,\qquad f(0,0)=10,\qquad
 f(1,\pm1)=-64,\qquad f(1,0)=108
\]
cross-check against the $\Delta_5$ denominator lane of
Theorem~\ref{thm:htbb-heegner-all-order}(ii).  The doubled values
belong to the scalar square
$\Phi_{10}^{\mathrm{un}}=\Delta_5^2$, not to the primitive
BKM root-supercharacter.
\end{remark}

\begin{remark}[Cross-volume invariance of the non-perturbative
completion]
\label{rem:thqg-np-cross-volume}
\index{cross-volume!non-perturbative completion}
Theorem~\ref{thm:thqg-np-d3-completion} is the Vol~II conditional
twisted-holographic member of the non-perturbative completion
triptych:
\begin{itemize}[leftmargin=1.8em]
\item Vol~I (\texttt{vol1/chapters/connections/bv\_brst.tex}
Theorem~$\mathrm{[thm:bvbrst-nonperturbative-completion]}$): Borel
resummation of the BV obstruction tower, with D$3$-instanton tail
from Costello~$2011$ Borel-plane analysis.
\item Vol~II (present conditional theorem plus
Theorem~\ref{thm:htbb-schwinger-humbert}): twisted-holographic
bulk--boundary--line partition function with Borcherds-product
automorphic packaging and Schwinger pair-creation at the Humbert
boundary, conditional on the resurgent hypotheses stated above.
\item Vol~III
(\texttt{vol3/chapters/theory/bkm\_denominator.tex}
Theorem~$\mathrm{[thm:bkm-phi10-denominator]}$): the Borcherds
product~\eqref{eq:thqg-np-borcherds-product} as the
Gritsenko--Nikulin denominator identity for $\mathfrak{g}_{\Delta_5}$,
indexing one factor per BKM root.
\end{itemize}
The three constructions agree on the Borcherds product term by term.
Interpreting that agreement as a non-perturbative partition function
of the holographic atom at $K3\times T^2$ is the conditional
resurgent reading of this section. This is the
non-perturbative analogue of the Heegner-pattern cross-volume
triangulation of Remark~\ref{rem:htbb-heegner-all-order-cross-volume}.
\end{remark}

\begin{remark}[Relation to Costello renormalisation]
\label{rem:thqg-np-costello-renorm}
\index{Costello renormalisation!Borel plane}
Costello~$2011$ Chapter~$4$~\cite[Ch.~4]{Costello2011} supplies the
BV-renormalisation framework in which saddle expansions and
one-loop determinants are defined. It does not, by itself, prove that
the twisted-holographic Heegner $\hbar$-series has D$3$ Borel
singularities with Borcherds residues. For twisted $11$D~SUGRA on
$\bR^3\times K3\times\bC^2$, the Euclidean D$3$-branes wrapping rigid
K$3$ curves (Proposition~\ref{prop:thqg-np-d3-amplitude}) and the
Quillen--Bismut determinants $\mathcal{Z}^{(1)}_{\det}(C)$ give the
candidate transseries sectors. Theorem~\ref{thm:thqg-np-borcherds-resummation}
proves the Borcherds product identity; Theorem~\ref{thm:thqg-np-d3-completion}
states the extra analytic hypotheses needed to read that product as a
Borel-resummed non-perturbative completion of the
Costello--Gaiotto--Paquette twisted holography
framework~\cite{CostelloGaiottoPaquette2018}.
\end{remark}
